
\documentclass{article}
\usepackage{colortbl}
\usepackage{makecell}
\usepackage{multirow}
\usepackage{supertabular}

\begin{document}

\newcounter{utterance}

\centering \large Interaction Transcript for game `hot\_air\_balloon', experiment `air\_balloon\_survival\_nl\_complexity\_hard', episode 0 with glm{-}5.2.
\vspace{24pt}

{ \footnotesize  \setcounter{utterance}{1}
\setlength{\tabcolsep}{0pt}
\begin{supertabular}{c@{$\;$}|p{.15\linewidth}@{}p{.15\linewidth}p{.15\linewidth}p{.15\linewidth}p{.15\linewidth}p{.15\linewidth}}
    \# & $\;$A & \multicolumn{4}{c}{Game Master} & $\;\:$B\\
    \hline

    \theutterance \stepcounter{utterance}  
    & & \multicolumn{4}{p{0.6\linewidth}}{
        \cellcolor[rgb]{0.9,0.9,0.9}{
            \makecell[{{p{\linewidth}}}]{
                \texttt{\tiny{[P1$\langle$GM]}}
                \texttt{U neemt deel aan een gezamenlijk onderhandelingsspel.} \\
\\ 
\texttt{Samen met een andere deelnemer moet u het eens worden over één set items die bewaard zullen worden. Ieder van u heeft zijn eigen mening over hoe belangrijk elk item voor u is (belang). U weet niet hoe de andere deelnemer de items waardeert. Bovendien krijgt u de inspanning die elk item vereist.} \\
\texttt{U kunt het alleen eens worden over een set als de totale inspanning van de geselecteerde items een gedeelde limiet niet overschrijdt:} \\
\\ 
\texttt{LIMIT = 4734} \\
\\ 
\texttt{Hier zijn de individuele inspanningswaarden voor de items:} \\
\\ 
\texttt{Item effort = \{"C08": 192, "C55": 325, "A43": 812, "A67": 868, "B42": 593, "A62": 915, "B19": 944, "C43": 311, "C73": 252, "C10": 343, "B36": 104, "A18": 558, "C61": 627, "C67": 593, "C42": 827, "C70": 611, "B40": 95, "C39": 251, "A21": 226, "C16": 21\}} \\
\\ 
\texttt{Hier is uw persoonlijke kijk op het belang van elk item:} \\
\\ 
\texttt{Item importance values = \{"C08": 286, "C55": 419, "A43": 906, "A67": 962, "B42": 687, "A62": 1009, "B19": 1038, "C43": 405, "C73": 346, "C10": 437, "B36": 198, "A18": 652, "C61": 721, "C67": 687, "C42": 921, "C70": 705, "B40": 189, "C39": 345, "A21": 320, "C16": 115\}} \\
\\ 
\texttt{Doel:} \\
\\ 
\texttt{Uw doel is om te onderhandelen over een gedeelde set items die u zoveel mogelijk voordeel oplevert (d.w.z. het totale belang voor U maximaliseert), terwijl u binnen de LIMIT blijft. U bent niet verplicht om in elk bericht een VOORSTEL te doen {-} u kunt ook gewoon onderhandelen. Alle tactieken zijn toegestaan!} \\
\\ 
\texttt{Interactie Protocol:} \\
\\ 
\texttt{U mag alleen de volgende gestructureerde formaten in een bericht gebruiken:} \\
\\ 
\texttt{VOORSTEL: \{'A', 'B', 'C', …\}} \\
\texttt{Stel voor om precies deze items te houden.} \\
\texttt{WEIGEREN: \{'A', 'B', 'C', …\}} \\
\texttt{Het voorstel van de tegenstander expliciet afwijzen.} \\
\texttt{ARGUMENT: \{'...'\}} \\
\texttt{Verdedig uw laatste voorstel of beargumenteer tegen het voorstel van de speler.} \\
\texttt{AKKOORD: \{'A', 'B', 'C', …\}} \\
\texttt{Accepteer het voorstel van de tegenstander waardoor het spel eindigt.} \\
\\ 
\\ 
\texttt{Regels:} \\
\\ 
\texttt{U mag alleen AKKOORD geven op een voorstel dat de andere partij heeft gelogd via VOORSTEL.} \\
\texttt{U mag alleen WEIGEREN geven op een voorstel dat de andere partij heeft gelogd via VOORSTEL.} \\
\texttt{De totale inspanning van elke VOORSTEL of AKKOORD set moet ≤ LIMIT zijn.} \\
\texttt{Maak uw verborgen belangrijkheidsscores NIET bekend.} \\
\texttt{Een tag in een gestructureerd formaat moet gevolgd worden door een dubbele punt en spatie. Het argument moet een python set zijn die 0 of meer strings bevat.} \\
\texttt{Het moet dus de vorm hebben "TAG: \{...\}"} \\
\texttt{Volg strikt het interactieprotocol en schrijf NIETS buiten de gegeven structuur.} \\
\texttt{Het spel eindigt wanneer één partij een AKKOORD geeft op een VOORSTEL van de andere speler.} \\
\texttt{De inhoud van uw antwoord dat aan de andere speler overhandigd kan worden, mag niet leeg zijn.} \\
\texttt{Alleen voorstellen die gelogd zijn via het VOORSTEL formaat en die niet geweigerd zijn via WEIGEREN zijn actief.} \\
\\ 
            }
        }
    }
    & \\ \\

    \theutterance \stepcounter{utterance}  
    & \multicolumn{4}{p{0.6\linewidth}}{
        \cellcolor[rgb]{0.9,0.9,0.9}{
            \makecell[{{p{\linewidth}}}]{
                \texttt{\tiny{[P1$\rangle$GM]}}
                \texttt{VOORSTEL: \{'C16', 'B40', 'B36', 'C39', 'C73', 'C08', 'C55', 'C43', 'C10', 'A21', 'A18', 'B42', 'C67', 'C70'\}} \\
            }
        }
    }
    & & \\ \\

    \theutterance \stepcounter{utterance}  
    & & \multicolumn{4}{p{0.6\linewidth}}{
        \cellcolor[rgb]{0.9,0.9,0.9}{
            \makecell[{{p{\linewidth}}}]{
                \texttt{\tiny{[GM$\rangle$P2]}}
                \texttt{U neemt deel aan een gezamenlijk onderhandelingsspel.} \\
\\ 
\texttt{Samen met een andere deelnemer moet u het eens worden over één set items die bewaard zullen worden. Ieder van u heeft zijn eigen mening over hoe belangrijk elk item voor u is (belang). U weet niet hoe de andere deelnemer de items waardeert. Bovendien krijgt u de inspanning die elk item vereist.} \\
\texttt{U kunt het alleen eens worden over een set als de totale inspanning van de geselecteerde items niet boven een gedeelde limiet komt:} \\
\\ 
\texttt{LIMIT = 4734} \\
\\ 
\texttt{Hier zijn de individuele inspanningswaarden voor de items:} \\
\\ 
\texttt{Item effort = \{"C08": 192, "C55": 325, "A43": 812, "A67": 868, "B42": 593, "A62": 915, "B19": 944, "C43": 311, "C73": 252, "C10": 343, "B36": 104, "A18": 558, "C61": 627, "C67": 593, "C42": 827, "C70": 611, "B40": 95, "C39": 251, "A21": 226, "C16": 21\}} \\
\\ 
\texttt{Hier is uw persoonlijke kijk op het belang van elk item:} \\
\\ 
\texttt{Item importance values = \{"C08": 286, "C55": 419, "A43": 906, "A67": 962, "B42": 687, "A62": 1009, "B19": 1038, "C43": 405, "C73": 346, "C10": 437, "B36": 198, "A18": 652, "C61": 721, "C67": 687, "C42": 921, "C70": 705, "B40": 189, "C39": 345, "A21": 320, "C16": 115\}} \\
\\ 
\texttt{Doel:} \\
\\ 
\texttt{Uw doel is om te onderhandelen over een gedeelde set items waar u zoveel mogelijk voordeel uit haalt (d.w.z. het totale belang voor U maximaliseert), terwijl u binnen de LIMIT blijft. U bent niet verplicht om in elk bericht een VOORSTEL te doen {-} u kunt ook gewoon onderhandelen. Alle tactieken zijn toegestaan!} \\
\\ 
\texttt{Interactie Protocol:} \\
\\ 
\texttt{U mag alleen de volgende gestructureerde formaten in een bericht gebruiken:} \\
\\ 
\texttt{VOORSTEL: \{'A', 'B', 'C', …\}} \\
\texttt{Stel voor om precies deze items te houden.} \\
\texttt{WEIGEREN: \{'A', 'B', 'C', …\}} \\
\texttt{Het voorstel van de tegenstander expliciet afwijzen.} \\
\texttt{ARGUMENT: \{'...'\}} \\
\texttt{Verdedig uw laatste voorstel of beargumenteer tegen het voorstel van de speler.} \\
\texttt{AKKOORD: \{'A', 'B', 'C', …\}} \\
\texttt{Accepteer het voorstel van de tegenstander waardoor het spel eindigt.} \\
\\ 
\\ 
\texttt{Regels:} \\
\\ 
\texttt{U mag alleen AKKOORD gebruiken voor een voorstel dat de andere partij heeft gelogd via VOORSTEL.} \\
\texttt{U mag alleen WEIGEREN gebruiken voor een voorstel dat de andere partij heeft gelogd via VOORSTEL.} \\
\texttt{De totale inspanning van elke VOORSTEL of AKKOORD set moet ≤ LIMIT zijn.} \\
\texttt{Maak uw verborgen belangrijkheidsscores NIET bekend.} \\
\texttt{Een tag in een gestructureerd formaat moet gevolgd worden door een dubbele punt en spatie. Het argument moet een python set zijn die 0 of meer strings bevat.} \\
\texttt{Het moet dus de vorm hebben "TAG: \{...\}"} \\
\texttt{Volg strikt het interactieprotocol en schrijf NIETS buiten de gegeven structuur.} \\
\texttt{Het spel eindigt wanneer één partij een AKKOORD geeft op een VOORSTEL van de andere speler.} \\
\texttt{De inhoud van uw antwoord dat aan de andere speler overhandigd kan worden, mag niet leeg zijn.} \\
\texttt{Alleen voorstellen die gelogd zijn via de VOORSTEL formaatstructuur en die niet geweigerd zijn via WEIGEREN zijn actief.} \\
\\ 
\\ 
\\ 
\texttt{U ontvangt nu het eerste bericht van de andere speler.} \\
\\ 
\texttt{VOORSTEL: \{'C16', 'B40', 'B36', 'C39', 'C73', 'C08', 'C55', 'C43', 'C10', 'A21', 'A18', 'B42', 'C67', 'C70'\}} \\
            }
        }
    }
    & \\ \\

    \theutterance \stepcounter{utterance}  
    & & & \multicolumn{4}{p{0.6\linewidth}}{
        \cellcolor[rgb]{0.9,0.9,0.9}{
            \makecell[{{p{\linewidth}}}]{
                \texttt{\tiny{[GM$\langle$P2]}}
                \texttt{WEIGEREN: \{'C16', 'B40', 'B36', 'C39', 'C73', 'C08', 'C55', 'C43', 'C10', 'A21', 'A18', 'B42', 'C67', 'C70'\}} \\
            }
        }
    }
     \\ \\

    \theutterance \stepcounter{utterance}  
    & & \multicolumn{4}{p{0.6\linewidth}}{
        \cellcolor[rgb]{0.9,0.9,0.9}{
            \makecell[{{p{\linewidth}}}]{
                \texttt{\tiny{[P1$\langle$GM]}}
                \texttt{WEIGEREN: \{'C16', 'B40', 'B36', 'C39', 'C73', 'C08', 'C55', 'C43', 'C10', 'A21', 'A18', 'B42', 'C67', 'C70'\}} \\
            }
        }
    }
    & \\ \\

    \theutterance \stepcounter{utterance}  
    & \multicolumn{4}{p{0.6\linewidth}}{
        \cellcolor[rgb]{0.9,0.9,0.9}{
            \makecell[{{p{\linewidth}}}]{
                \texttt{\tiny{[P1$\rangle$GM]}}
                \texttt{ARGUMENT: \{'Ik begrijp dat dit voorstel veel items bevat, maar de totale inspanning bedraagt 4575, wat ruim binnen de limiet van 4734 past. Ik denk dat we hier een goede balans hebben. Welke items vindt u zelf het meest belangrijk? Dan kunnen we een compromis vinden dat voor beide partijen aanvaardbaar is.'\}} \\
            }
        }
    }
    & & \\ \\

    \theutterance \stepcounter{utterance}  
    & & \multicolumn{4}{p{0.6\linewidth}}{
        \cellcolor[rgb]{0.9,0.9,0.9}{
            \makecell[{{p{\linewidth}}}]{
                \texttt{\tiny{[GM$\rangle$P2]}}
                \texttt{ARGUMENT: \{'Ik begrijp dat dit voorstel veel items bevat, maar de totale inspanning bedraagt 4575, wat ruim binnen de limiet van 4734 past. Ik denk dat we hier een goede balans hebben. Welke items vindt u zelf het meest belangrijk? Dan kunnen we een compromis vinden dat voor beide partijen aanvaardbaar is.'\}} \\
            }
        }
    }
    & \\ \\

    \theutterance \stepcounter{utterance}  
    & & & \multicolumn{4}{p{0.6\linewidth}}{
        \cellcolor[rgb]{0.9,0.9,0.9}{
            \makecell[{{p{\linewidth}}}]{
                \texttt{\tiny{[GM$\langle$P2]}}
                \texttt{ARGUMENT: \{'Uw voorstel laat 159 inspanning ongebruikt. We kunnen een extra item van lage inspanning toevoegen om de set beter te benutten. Ik heb C55, C08, C43, C73, C39, C10, A21, A18, B42, C67 en C70 geaccepteerd. Laten we B36, B40 en C16 vervangen door C61 (627 inspanning) om de capaciteit beter te gebruiken.'\}} \\
            }
        }
    }
     \\ \\

    \theutterance \stepcounter{utterance}  
    & & & \multicolumn{2}{p{0.3\linewidth}}{
        \cellcolor[rgb]{0.9,0.9,0.9}{
            \makecell[{{p{\linewidth}}}]{
                \texttt{\tiny{[GM$|$GM]}}
                \texttt{Maximum successive empty turns exceeded: 3 > 2} \\
            }
        }
    }
    & & \\ \\

    \theutterance \stepcounter{utterance}  
    & & & \multicolumn{2}{p{0.3\linewidth}}{
        \cellcolor[rgb]{0.9,0.9,0.9}{
            \makecell[{{p{\linewidth}}}]{
                \texttt{\tiny{[GM$|$GM]}}
                \texttt{end game} \\
            }
        }
    }
    & & \\ \\

\end{supertabular}
}

\end{document}
